\subsection{A coefficient-two exceptional curve over an ordinary-triple germ}

\begin{statement}{P4.39}{Proved}
Let $S=(f=0)$ be a reduced hypersurface germ in a smooth
threefold, with one-dimensional reduced critical locus $C$.
Suppose that every generic transverse germ along $C$ is an
ordinary triple point. On the minimal resolution $T$ of its
normalization, let $D$ be the integral total pole divisor
defined by the residue of an ambient volume form divided by
$f$. No exceptional curve of positive genus has coefficient
two in $D$.

In particular, the contracted horizontal section of a
positive-genus ruled resolution in Part~4.2 cannot lie over
a point where all local critical-curve branches have generic
type $D_4$. This assertion does not replace that local
hypothesis by a mixture of transverse types.
\end{statement}
\begin{proof}
We blow up the ambient origin, write $S_1$ for the strict
transform and $N_1$ for its normalization, and take the
minimal resolution $T_1$ of $N_1$. The proper map $N_1\to S$
lifts to the normalization $N$ by its universal property.
Thus $T_1\to N$ is a resolution and dominates $T$ by a
sequence of point blowups. The strict transform of a proposed
positive-genus exceptional curve $H$ remains present and has
the same pole coefficient two.

Put $d=\operatorname{mult}_0f$, and denote the ambient
exceptional divisor by $E$. Actual adjunction gives
\[
 \omega_{S_1}=\pi^*\omega_S\otimes
        \mathcal O((2-d)E)|_{S_1}.
\]
Let $D_1$ be the total conductor divisor on $T_1$ relative
to $\omega_{S_1}$, and let $e$ be the Cartier pullback of
$E$. With the specified residue comparison, the total pole
of the original germ on $T_1$ is
\[
                    D_1+(d-2)e.                       \tag{4.39a}
\]
The divisor $D_1$ is effective. Its strict conductor
coefficients are positive integers. For exceptional curves
of the minimal resolution of $N_1$, the negative-matrix
argument of P4R.6 applies locally: $K_{T_1}F\ge0$, since
there is no smooth rational exceptional $(-1)$-curve;
adding the strict conductor retains nonnegative intersection.
The inverse exceptional intersection matrix has nonpositive
entries, so all exceptional coefficients are nonnegative.

If $H$ is exceptional over $N_1$, its coefficient in $D_1$
is strictly positive. A zero coefficient would give
\[
 0=(K_{T_1}+D_1)H
   =2g(H)-2-H^2+D_1H>0.
\]
If that coefficient is one, adjunction gives
\[
 0=2g(H)-2+(D_1-H)H.
\]
Hence $H$ is elliptic and is disjoint from every other
positive component of $D_1$. A zero-coefficient exceptional
curve cannot meet a positive one, by its own intersection
equation and $K_{T_1}F\ge0$. The connected exceptional
fibre is consequently $H$ alone and meets no strict conductor.
The entire germ of $S_1$ at its image is normal with an
isolated simple elliptic singularity. Indeed, any additional
hypersurface branch would intersect this branch in a curve
and would give strict conductor meeting $H$.

Set $n=\operatorname{ord}_H(e)>0$ and
$a=\operatorname{coeff}_H(D_1)\ge0$. Formula (4.39a)
gives the exact identity
\[
                       2=(d-2)n+a.                    \tag{4.39b}
\]
Thus $d\ge5$ is impossible. For $d=4$ we have
$n=1,a=0$, and $H$ is not exceptional over $N_1$.
At its generic point the normalization is an isomorphism,
so $H$ is birational to a reduced component of the projective
quartic tangent cone. Every second derivative of $f$ vanishes
on $C$. Its leading term along any analytic arc of $C$
therefore makes that quartic have a point of multiplicity
at least three. An irreducible quartic has arithmetic genus
three and loses at least three there. A smooth elliptic cubic
component together with the remaining line could have
multiplicity at most two at the point. All components of
degree at most two are rational. Thus every reduced tangent
component is rational, excluding $H$ for $d=4$.

We are left with $d=3$ and
\[
                   (n,a)=(2,0)\quad\hbox{or}\quad(1,1).
\]
Choose a coordinate $u$ with nonzero pure cubic initial
coefficient, apply Weierstrass preparation, and depress the
monic cubic. Multiplication by the preparation unit does not
affect Hessian vanishing on $C$, because $f=df=0$ there.
We obtain
\[
 f=\text{unit}\cdot\bigl(u^3+b(x,z)u+c(x,z)\bigr).
\]
The second $u$ derivative is $6u$, so $C$ lies in the smooth
surface $u=0$. Write its reduced ideal as $(u,h(x,z))$.
The complete-intersection cube is unmixed and equals its
symbolic cube. Generic order three along $C$ and coefficient
comparison therefore give the exact expression
\[
 f=\text{unit}\cdot
       \bigl(u^3+h^2B(x,z)u+h^3C(x,z)\bigr).           \tag{4.39c}
\]
If $C$ is singular, $\operatorname{ord}h\ge2$ and the
tangent cone is the triple plane $u^3$. The degree-four
term restricts to zero on this plane. After the ambient
blowup the entire reduced exceptional line is critical,
so none of its points is an isolated normal singularity.
If $C$ is smooth, take $x=h$. Every homogeneous term $f_j$
then has degree at least three in $u,x$. The cubic tangent
cone is a binary cubic in these two variables, hence three
lines, counted with multiplicities, through the tangent
direction of $C$.

First suppose that $H$ dominates a tangent-cone line.
For $(n,a)=(2,0)$ the normalization is generically an
isomorphism, and $H$ is birational to that line. For
$(n,a)=(1,1)$, a transverse normalization branch has conductor
exponent one. It is smooth and meets exactly one other smooth
branch transversely: the intrinsic conductor exponent of a
singular plane branch is at least two, and each other branch
contributes a positive intersection multiplicity. Thus the
generic transverse germ of $S_1$ is a node. A multiplicity-one
tangent line is generically smooth and cannot have $a=1$.
For a multiplicity-three line, the two geometric branches of
that node have exceptional orders whose weighted sum is three.
A connected double cover would give equal orders and an even
sum. Thus the order-one branch has degree one over the line
and is rational.

The remaining possibility is a double line together with a
distinct simple line. Here $C$ is smooth by (4.39c). Make a
linear change so that $f_3=xu^2$. On the $x$ chart of the
blowup put $u=xv,z=xw$. The strict equation is
\[
 v^2+x f_4(1,v,w)+x^2f_5(1,v,w)+\cdots.
\]
An ordinary double curve on $x=v=0$ requires
$f_4(1,0,w)=0$ identically. Its transverse quadratic form is
\[
 v^2+xvB(w)+x^2C(w),
 \qquad \deg B\le1,\quad\deg C\le2.                  \tag{4.39d}
\]
The degree bounds follow from the degree-at-least-three
condition in $u,x$ on every $f_j$: here
$B=\partial_v f_4(1,0,w)$ and $C=f_5(1,0,w)$.
The normalization over the exceptional line adjoins a root
of $q^2+B(w)q+C(w)$. Its discriminant has degree at most
two. The resulting quadratic extension of $\mathbb C(w)$
is branched at at most two points, including infinity.
Riemann--Hurwitz makes its compact normalization rational;
the components of a split cover are rational as well.

Finally suppose that $H$ is exceptional over a point of $N_1$.
We must have $n=a=1$, and the first part of the proof would
make $S_1$ normal with an isolated simple elliptic singularity
there. This cannot happen. For singular $C$ the whole
exceptional line is critical, as already proved. For smooth
$C$ its strict transform passes through the common point of
the tangent lines, so that point is not an isolated normal
singularity. Away from it, at a line of multiplicity
$e=2$ or $3$, divide by the nonzero factor from any other
line. The strict equation has the form
\[
                         v^e+xG(v,w,x),
\]
where $G(0,w,0)$ is affine linear in $w$, by the same
normal-degree bound. If this function is nonzero at a point,
the surface is smooth. If it is identically zero, the entire
line is critical. At a simple zero, $G_w$ is a unit and the
exact coordinate change $W=G(v,w,x)$ gives $v^e+xW$, a
rational double point of type $A_{e-1}$. A multiplicity-one
line is smooth away from intersections. There is therefore
no isolated simple elliptic singularity. All possibilities
for $H$ have been excluded.
\end{proof}

The positive-integer solutions of (4.39b), the monomial degree
bounds in (4.39d), and the forced fourth-order factor of the
homogeneous quadratic discriminant are checked in
\texttt{code/part4/d4\_pole.py}. Its saved output reports PASS.
